% theory_endogenous_premium.tex
% Central theoretical contribution: Endogenous NAV Premium Theory
%
% Notation follows plan.tex: S_t (BTC price), H_t (holdings), pi_t (log-premium),
% kappa (mean-reversion speed), theta (long-run mean), sigma_S, sigma_pi, etc.

\newtheorem{theorem}{Theorem}[section]
\newtheorem{proposition}[theorem]{Proposition}
\newtheorem{lemma}[theorem]{Lemma}
\newtheorem{corollary}[theorem]{Corollary}
\newtheorem{definition}[theorem]{Definition}
\newtheorem{assumption}[theorem]{Assumption}
\newtheorem{remark}[theorem]{Remark}

%=============================================================================
\section{Endogenous Premium Theory}
\label{sec:endogenous_premium}
%=============================================================================

The baseline model of Section~2 treats the log-premium $\pi_t$ as an
exogenous Ornstein--Uhlenbeck process mean-reverting to a constant $\theta$.
This section derives $\theta$ endogenously: the equilibrium premium
$\pi^\star$ equals the market value of a \emph{perpetual accumulation option}
--- the option to issue equity at a premium and purchase BTC, thereby
increasing BTC-per-share for existing holders.

%-----------------------------------------------------------------------------
\subsection{Setup and Assumptions}
\label{sec:endo_setup}
%-----------------------------------------------------------------------------

We retain the filtered probability space
$(\Omega,\mathcal{F},(\mathcal{F}_t)_{t\ge 0},\mathbb{P})$ and the state
variables from the baseline model:
\begin{itemize}
  \item $S_t$: BTC spot price;
  \item $H_t$: BTC holdings (in BTC);
  \item $D_t$: outstanding debt (face value);
  \item $N_t^{\mathrm{sh}}$: shares outstanding;
  \item $B_t = H_t / N_t^{\mathrm{sh}}$: BTC per share;
  \item $\mathrm{NAV}_t = (H_t S_t - D_t)^+$: net asset value;
  \item $E_t = e^{\pi_t}\,\mathrm{NAV}_t$: equity market capitalization.
\end{itemize}

We collect the observable state into the vector
\begin{equation}
  \label{eq:state_vector}
  \Phi_t = (S_t,\, H_t,\, D_t,\, N_t^{\mathrm{sh}},\, \sigma_S,\, \kappa,\, \alpha),
\end{equation}
where $\sigma_S$ is BTC volatility and $\alpha$ and $\kappa$ are behavioral
parameters defined below.

\begin{assumption}[Operating regime]
  \label{ass:operating_regime}
  Throughout this section we assume $A_t = H_t S_t > D_t$ so that
  $\mathrm{NAV}_t > 0$ and the log-premium $\pi_t = \log(E_t/\mathrm{NAV}_t)$
  is well-defined.  The distress regime $A_t \le D_t$ is treated separately
  in Section~\ref{sec:death_spiral}.
\end{assumption}

\begin{assumption}[Equity issuance technology]
  \label{ass:issuance}
  The firm can issue new shares at-the-market at the prevailing equity price
  $P_t$ and immediately deploy the proceeds to purchase BTC at price $S_t$.
  Issuance is frictionless up to a market-impact cost parameterized by $\eta \ge 0$.
\end{assumption}

%-----------------------------------------------------------------------------
\subsection{Endogenous Premium System}
\label{sec:coupled_system}
%-----------------------------------------------------------------------------

We replace the exogenous OU premium with a \emph{coupled} system of SDEs.

\begin{definition}[Endogenous premium system]
  \label{def:coupled_sdes}
  The joint dynamics of $(S_t, \pi_t, H_t)$ under $\mathbb{P}$ are:
  %
  \begin{align}
    \frac{dS_t}{S_{t^-}} &= \mu_S\,dt + \sigma_S\,dW_t^S
      + \eta\,\frac{dH_t}{H_{t^-}},
    \label{eq:btc_impact} \\[6pt]
    d\pi_t &= \kappa\bigl(\pi^\star(\Phi_t) - \pi_t\bigr)\,dt
      + \sigma_\pi\,dW_t^\pi,
    \label{eq:premium_endo} \\[6pt]
    dH_t &= \alpha\,\pi_t^+\,dt + dH_t^{\mathrm{jump}},
    \label{eq:holdings_endo}
  \end{align}
  %
  where:
  \begin{itemize}
    \item $\mu_S$ is the BTC drift under $\mathbb{P}$;
    \item $\eta \ge 0$ is the market-impact elasticity of BTC price to MSTR
          accumulation;
    \item $\pi^\star(\Phi_t)$ is the \emph{endogenous fair premium} derived
          below;
    \item $\alpha > 0$ governs continuous accumulation intensity;
    \item $\pi_t^+ = \max(\pi_t,0)$;
    \item $dH_t^{\mathrm{jump}} = \sum_k \Delta H_{\tau_k}\,\delta(t - \tau_k)$
          captures discrete financing events with Poisson arrival rate
          $\lambda_M$;
    \item $\mathrm{corr}(dW_t^S, dW_t^\pi) = \rho$.
  \end{itemize}
\end{definition}

\begin{remark}
  Equation~\eqref{eq:premium_endo} nests the baseline model: setting
  $\pi^\star(\Phi_t) \equiv \theta$ (constant) recovers the exogenous OU
  process of plan.tex, Section~2.3.
\end{remark}

The key novelty is that the long-run target $\pi^\star$ is itself a function
of the state $\Phi_t$, creating a \emph{feedback loop}:
\begin{equation}
  \label{eq:feedback_loop}
  \pi_t \uparrow \;\Longrightarrow\;
  H_t \uparrow \;\Longrightarrow\;
  B_t \uparrow \;\Longrightarrow\;
  \pi^\star \uparrow \;\Longrightarrow\;
  \pi_t \uparrow \;\cdots
\end{equation}
This is the \emph{reflexive flywheel} that distinguishes the MSTR structure
from a passive BTC fund.

%=============================================================================
\subsection{The Accumulation Option and Fair Premium}
\label{sec:fair_premium}
%=============================================================================

%-----------------------------------------------------------------------------
\subsubsection{Accretive Issuance: Single-Round Mechanics}
\label{sec:single_round}
%-----------------------------------------------------------------------------

\begin{definition}[Accretion from a single issuance]
  \label{def:single_accretion}
  Suppose the firm has $H$ BTC, $N$ shares, and trades at premium $\pi > 0$.
  It issues $\delta N$ new shares at the market price
  $P = e^\pi \cdot (HS - D)/N$, raising proceeds
  $\delta N \cdot P = e^\pi \cdot \delta N \cdot (HS - D)/N$, and purchases
  $\delta H$ BTC at price $S$.  Then:
  %
  \begin{equation}
    \label{eq:delta_H}
    \delta H = \frac{e^\pi (HS - D)}{NS}\,\delta N,
  \end{equation}
  %
  and the change in BTC-per-share is
  %
  \begin{equation}
    \label{eq:accretion_single}
    \frac{\Delta B}{B}
    = \frac{H + \delta H}{N + \delta N}\cdot\frac{N}{H} - 1
    = \frac{(e^\pi - 1)\,(1 - D/(HS))}{1 + \delta N / N}\cdot
      \frac{\delta N}{N}.
  \end{equation}
\end{definition}

For small issuance ($\delta N / N \ll 1$) and ignoring debt for clarity:
\begin{equation}
  \label{eq:accretion_approx}
  \frac{\Delta B}{B} \approx (e^\pi - 1)\,\frac{\delta N}{N}.
\end{equation}
Each issuance round at premium $\pi > 0$ is BTC-per-share accretive.  At
$\pi < 0$ (discount to NAV), issuance would be dilutive, so the firm
optimally refrains: this is the option-like payoff.

%-----------------------------------------------------------------------------
\subsubsection{Perpetual Accumulation Option Value}
\label{sec:vacc}
%-----------------------------------------------------------------------------

\begin{definition}[Accumulation option value]
  \label{def:vacc}
  Define the \emph{accumulation option value} as the present value (under
  $\mathbb{P}$) of all future BTC-per-share accretion from optimally timed
  issuance:
  %
  \begin{equation}
    \label{eq:vacc_def}
    V_{\mathrm{acc}}(\Phi_t)
    = \mathrm{NAV}_t \cdot
      \mathbb{E}^{\mathbb{P}}\!\left[
        \int_t^\infty e^{-r(u-t)}\,
        g(\pi_u,\Phi_u)\,du
        \;\Big|\;\mathcal{F}_t
      \right],
  \end{equation}
  %
  where $r > 0$ is the discount rate and the \emph{accretion rate function}
  is
  \begin{equation}
    \label{eq:accretion_rate}
    g(\pi,\Phi)
    = \alpha_N\,(e^\pi - 1)^+\,\biggl(1 - \frac{D}{HS}\biggr)^+,
  \end{equation}
  with $\alpha_N > 0$ representing the issuance intensity (shares issued per
  unit time as a fraction of $N$).
\end{definition}

The term $(e^\pi - 1)^+$ captures the option-like feature: the firm issues
only when $\pi > 0$.  The factor $(1 - D/(HS))^+$ adjusts for leverage.

\begin{proposition}[Fair premium]
  \label{prop:fair_premium}
  Under the endogenous premium framework, the \emph{fair (equilibrium)
  log-premium} satisfies
  \begin{equation}
    \label{eq:pi_star_def}
    \boxed{
      \pi^\star(\Phi_t)
      = \log\!\left(1 + \frac{V_{\mathrm{acc}}(\Phi_t)}{\mathrm{NAV}_t}\right).
    }
  \end{equation}
\end{proposition}

\begin{proof}
  In equilibrium, equity value equals NAV plus the value of the accumulation
  option:
  \[
    E_t = \mathrm{NAV}_t + V_{\mathrm{acc}}(\Phi_t).
  \]
  Since $E_t = e^{\pi_t}\,\mathrm{NAV}_t$ by definition, at equilibrium
  $\pi_t = \pi^\star$ we obtain
  \[
    e^{\pi^\star}\,\mathrm{NAV}_t = \mathrm{NAV}_t + V_{\mathrm{acc}},
  \]
  which gives~\eqref{eq:pi_star_def}.
\end{proof}

\begin{remark}
  Equation~\eqref{eq:pi_star_def} provides economic content to the premium:
  $\pi^\star > 0$ if and only if the accumulation option has positive value,
  which requires $\mathbb{P}(\pi_u > 0 \text{ for some } u > t) > 0$.
  A passive BTC holding company with no issuance capability ($\alpha_N = 0$)
  has $V_{\mathrm{acc}} = 0$ and thus $\pi^\star = 0$: it should trade at NAV.
\end{remark}

%-----------------------------------------------------------------------------
\subsubsection{Steady-State Closed-Form Approximation}
\label{sec:steady_state}
%-----------------------------------------------------------------------------

To obtain a tractable expression, we approximate the premium dynamics at
steady state and linearize the accretion payoff.

\begin{assumption}[Steady-state approximation]
  \label{ass:steady_state}
  At steady state, the premium $\pi_t$ fluctuates around $\pi^\star$ as a
  stationary OU process with speed $\kappa$ and volatility $\sigma_\pi$.
  The state variables $(S_t, H_t, D_t)$ evolve slowly relative to premium
  mean-reversion, so we treat them as locally constant.
\end{assumption}

Under Assumption~\ref{ass:steady_state}, the premium has stationary
distribution $\pi \sim \mathcal{N}(\pi^\star, \sigma_\pi^2/(2\kappa))$.
Define $\sigma_\pi^{\mathrm{ss}} = \sigma_\pi / \sqrt{2\kappa}$.

\begin{proposition}[Steady-state fair premium]
  \label{prop:steady_state}
  Under Assumption~\ref{ass:steady_state}, with leverage ratio
  $\ell = D/(HS)$ and $L = (1-\ell)^+$, the fair premium satisfies the
  fixed-point equation
  \begin{equation}
    \label{eq:pi_star_fp}
    \pi^\star
    = \log\!\left(
        1 + \frac{\alpha_N L}{r}\,
        \mathbb{E}\bigl[(e^{\pi} - 1)^+\bigr]
      \right),
  \end{equation}
  where $\pi \sim \mathcal{N}(\pi^\star,\, (\sigma_\pi^{\mathrm{ss}})^2)$.
  Evaluating the expectation yields
  \begin{equation}
    \label{eq:pi_star_closed}
    \boxed{
      \pi^\star
      = \log\!\left(
          1 + \frac{\alpha_N L}{r}\,\Bigl[
            e^{\pi^\star + (\sigma_\pi^{\mathrm{ss}})^2/2}\,
            \Phi\!\left(\frac{\pi^\star}{\sigma_\pi^{\mathrm{ss}}}
              + \sigma_\pi^{\mathrm{ss}}\right)
            - \Phi\!\left(\frac{\pi^\star}{\sigma_\pi^{\mathrm{ss}}}\right)
          \Bigr]
        \right),
    }
  \end{equation}
  where $\Phi(\cdot)$ is the standard normal CDF.
\end{proposition}

\begin{proof}
  At steady state, future accretion discounted at rate $r$ gives
  \[
    \frac{V_{\mathrm{acc}}}{\mathrm{NAV}}
    = \frac{\alpha_N L}{r}\,\mathbb{E}[(e^\pi - 1)^+].
  \]
  The expectation $\mathbb{E}[(e^\pi - 1)^+]$ with
  $\pi \sim \mathcal{N}(m, v^2)$ is the standard Black--Scholes call-option
  formula evaluated at strike $K=1$ on the underlying $e^\pi$:
  \[
    \mathbb{E}[(e^\pi - 1)^+]
    = e^{m + v^2/2}\,\Phi\!\left(\frac{m}{v} + v\right)
      - \Phi\!\left(\frac{m}{v}\right),
  \]
  with $m = \pi^\star$ and $v = \sigma_\pi^{\mathrm{ss}}$.
  Substituting into~\eqref{eq:pi_star_def} yields~\eqref{eq:pi_star_closed}.
\end{proof}

\begin{remark}
  Equation~\eqref{eq:pi_star_closed} is an implicit (fixed-point) equation in
  $\pi^\star$ because the right-hand side depends on $\pi^\star$.
  Existence and multiplicity of solutions are analyzed in
  Section~\ref{sec:multiple_equilibria}.
\end{remark}

%-----------------------------------------------------------------------------
\subsubsection{Comparative Statics}
\label{sec:comparative_statics}
%-----------------------------------------------------------------------------

\begin{proposition}[Comparative statics of the fair premium]
  \label{prop:comp_statics}
  Define $F(\pi^\star;\,\sigma_S,\kappa,\alpha_N,\ell,r)$ to be the
  right-hand side of~\eqref{eq:pi_star_closed}.  At any interior fixed point
  $\pi^\star > 0$ with $\partial F/\partial \pi^\star < 1$ (stable
  equilibrium):
  %
  \begin{enumerate}
    \item \textbf{BTC volatility:}
      $\displaystyle\frac{\partial \pi^\star}{\partial \sigma_S} > 0.$
      Higher BTC volatility increases premium volatility
      ($\sigma_\pi$ is increasing in $\sigma_S$ through the correlation
      $\rho$), raising the option value of future accretive issuance.

    \item \textbf{Mean-reversion speed:}
      $\displaystyle\frac{\partial \pi^\star}{\partial \kappa} < 0.$
      Faster mean-reversion reduces the stationary variance
      $\sigma_\pi^2/(2\kappa)$ and hence the option value, lowering the
      equilibrium premium.

    \item \textbf{Issuance intensity:}
      $\displaystyle\frac{\partial \pi^\star}{\partial \alpha_N} > 0.$
      Greater ability to issue equity increases the accumulation option value.
      This is the IBGR channel.

    \item \textbf{Leverage:}
      $\displaystyle\frac{\partial \pi^\star}{\partial \ell} < 0$ for
      $\ell < 1.$
      Higher leverage ($\ell = D/(HS) \uparrow$) reduces $L = 1 - \ell$ and
      hence the accretion efficiency per issuance round.

    \item \textbf{Discount rate:}
      $\displaystyle\frac{\partial \pi^\star}{\partial r} < 0.$
      Higher discount rates reduce the present value of future accretion.
  \end{enumerate}
\end{proposition}

\begin{proof}
  By the implicit function theorem, at a fixed point where
  $\pi^\star = F(\pi^\star; \mathbf{p})$ with $F_{\pi^\star} < 1$,
  \[
    \frac{\partial \pi^\star}{\partial p_i}
    = \frac{\partial F / \partial p_i}{1 - \partial F / \partial \pi^\star}.
  \]
  The denominator is positive at a stable equilibrium.  The signs of the
  numerators follow from direct differentiation of~\eqref{eq:pi_star_closed}:

  (1) $\sigma_\pi^{\mathrm{ss}} = \sigma_\pi/\sqrt{2\kappa}$ is increasing in
  $\sigma_\pi$ (and hence in $\sigma_S$ via $\rho$).  The Black--Scholes call
  value is increasing in volatility (positive vega), so $\partial F/\partial
  \sigma_S > 0$.

  (2) $\sigma_\pi^{\mathrm{ss}}$ is decreasing in $\kappa$, so the option
  value and hence $F$ decrease: $\partial F/\partial\kappa < 0$.

  (3) $F$ is manifestly increasing in $\alpha_N$ (multiplicative in the option
  value).

  (4) $L = (1-\ell)^+$ is decreasing in $\ell$ for $\ell < 1$, and $F$ is
  increasing in $L$.

  (5) $F$ is decreasing in $r$ since $r$ appears in the denominator.
\end{proof}

\begin{corollary}[Conditions for positive vs.\ negative fair premium]
  \label{cor:sign_pistar}
  The fair premium satisfies:
  \begin{enumerate}
    \item $\pi^\star > 0$ if and only if
      $\alpha_N > 0$, $L > 0$ (i.e., $\ell < 1$), and the premium
      volatility $\sigma_\pi^{\mathrm{ss}} > 0$.
    \item $\pi^\star = 0$ if the firm has no issuance capability
      ($\alpha_N = 0$), is at maximum leverage ($\ell \ge 1$), or if
      premium volatility vanishes.
    \item $\pi^\star < 0$ does not arise from the accumulation option alone;
      negative premia require additional factors (e.g., agency costs,
      management fee drag, illiquidity discount) not modeled here.
  \end{enumerate}
\end{corollary}

%=============================================================================
\subsection{Reflexivity and Multiple Equilibria}
\label{sec:multiple_equilibria}
%=============================================================================

The defining feature of the MSTR structure is a \emph{reflexive feedback loop}:
the premium determines accumulation, which determines the premium.  This
section formalizes the loop and characterizes the resulting equilibrium
multiplicity.

%-----------------------------------------------------------------------------
\subsubsection{Feedback Gain}
\label{sec:feedback_gain}
%-----------------------------------------------------------------------------

\begin{definition}[Feedback gain]
  \label{def:feedback_gain}
  The \emph{feedback gain} at equilibrium $\pi^\star$ is
  \begin{equation}
    \label{eq:feedback_gain}
    G(\pi^\star)
    = \frac{\partial F}{\partial \pi^\star}
    = \frac{\partial \pi^\star_{\mathrm{RHS}}}{\partial \pi^\star},
  \end{equation}
  where $F$ is the right-hand side of~\eqref{eq:pi_star_closed}.
  Equivalently, $G$ measures the marginal increase in the fair premium from a
  unit increase in the current premium level.
\end{definition}

\begin{proposition}[Effective mean-reversion]
  \label{prop:effective_kappa}
  Near a stable equilibrium $\pi^\star$ with feedback gain $G < 1$, the
  premium dynamics~\eqref{eq:premium_endo} exhibit an \emph{effective
  mean-reversion speed}
  \begin{equation}
    \label{eq:kappa_eff}
    \boxed{
      \kappa_{\mathrm{eff}} = \kappa\,(1 - G).
    }
  \end{equation}
  The feedback loop slows mean-reversion by the factor $(1-G)$.
\end{proposition}

\begin{proof}
  Linearize $\pi^\star(\Phi_t)$ around a steady state.  When $\pi_t$ increases
  by $\delta\pi$, the equilibrium target shifts by $G\,\delta\pi$.  The
  restoring drift becomes
  \[
    \kappa\bigl(\pi^\star + G\,\delta\pi - (\pi^\star + \delta\pi)\bigr)
    = -\kappa(1 - G)\,\delta\pi,
  \]
  yielding effective speed $\kappa_{\mathrm{eff}} = \kappa(1-G)$.
\end{proof}

\begin{remark}
  Three regimes emerge:
  \begin{itemize}
    \item $G < 1$: stable equilibrium; premium mean-reverts more slowly than
          in the exogenous model.
    \item $G = 1$: critical point; mean-reversion vanishes and the premium
          becomes a random walk locally.
    \item $G > 1$: unstable; small perturbations are self-amplifying (the
          ``flywheel'' or ``death spiral'').
  \end{itemize}
\end{remark}

%-----------------------------------------------------------------------------
\subsubsection{Fixed-Point Analysis and Equilibrium Multiplicity}
\label{sec:fixed_point}
%-----------------------------------------------------------------------------

\begin{theorem}[Existence and multiplicity of equilibria]
  \label{thm:three_equilibria}
  Consider the fixed-point equation $\pi^\star = F(\pi^\star)$
  from~\eqref{eq:pi_star_closed} with parameters
  $(\alpha_N, L, r, \sigma_\pi^{\mathrm{ss}})$ all strictly positive.
  Define the \emph{feedback strength}
  \begin{equation}
    \label{eq:feedback_strength}
    \Gamma = \frac{\alpha_N L}{r}.
  \end{equation}
  Then:
  %
  \begin{enumerate}
    \item \textbf{(Weak feedback)} If $\Gamma$ is sufficiently small, there
      exists a \emph{unique} fixed point $\pi^\star_{\mathrm{low}} > 0$,
      which is stable ($G < 1$).  This is the \emph{normal premium regime}.

    \item \textbf{(Strong feedback)} If $\Gamma$ is sufficiently large,
      there exist \emph{three} fixed points:
      \begin{itemize}
        \item $\pi^\star_{\mathrm{low}}$: the \emph{low/spiral equilibrium}
              (stable, $G < 1$);
        \item $\pi^\star_{\mathrm{mid}}$: the \emph{tipping point}
              (unstable, $G > 1$);
        \item $\pi^\star_{\mathrm{high}}$: the \emph{flywheel equilibrium}
              (stable, $G < 1$).
      \end{itemize}

    \item \textbf{(Bifurcation)} The transition between one and three
      equilibria occurs at a critical $\Gamma^*$ where $\pi^\star_{\mathrm{mid}}$
      and $\pi^\star_{\mathrm{high}}$ collide in a \emph{saddle-node
      bifurcation}.
  \end{enumerate}
\end{theorem}

\begin{proof}
  The function $F(\pi^\star)$ from~\eqref{eq:pi_star_closed} satisfies the
  following properties:

  \emph{Step 1: Boundary behavior.}
  $F(0) > 0$ (the option value is strictly positive when
  $\sigma_\pi^{\mathrm{ss}} > 0$, since $\mathbb{E}[(e^\pi - 1)^+] > 0$ for
  any non-degenerate distribution).  As $\pi^\star \to \infty$,
  $F(\pi^\star) \sim \log(1 + \Gamma e^{\pi^\star}) \sim \pi^\star +
  \log\Gamma$, so $F(\pi^\star) - \pi^\star \to \log\Gamma$.
  However, the growth rate of $F$ in $\pi^\star$ eventually exceeds 1 for
  intermediate values of $\pi^\star$ when $\Gamma$ is large, due to the
  exponential term $e^{\pi^\star + (\sigma_\pi^{\mathrm{ss}})^2/2}$ in the
  option value.

  \emph{Step 2: Shape of $F$.}
  Define $h(\pi^\star) = F(\pi^\star) - \pi^\star$.  We have $h(0) > 0$.
  For small $\Gamma$, $F'(\pi^\star) < 1$ everywhere, so $h$ is strictly
  decreasing and has a unique zero.

  For large $\Gamma$, $F'(\pi^\star)$ can exceed 1 over an intermediate
  range (the exponential growth of the option value dominates the logarithmic
  compression), creating a region where $h$ is increasing.  This produces
  the characteristic S-shaped curve with three crossings of $h = 0$.

  \emph{Step 3: Saddle-node bifurcation.}
  At the critical $\Gamma^*$, the function $h$ has a degenerate zero where
  $h(\pi^\star_c) = 0$ and $h'(\pi^\star_c) = 0$ simultaneously, i.e.,
  $F(\pi^\star_c) = \pi^\star_c$ and $F'(\pi^\star_c) = 1$.  This is the
  standard saddle-node (fold) bifurcation.  For $\Gamma$ slightly above
  $\Gamma^*$, the degenerate zero splits into two (the tipping point and the
  flywheel equilibrium), while the low equilibrium persists.

  \emph{Step 4: Stability.}
  At any fixed point, the linearized dynamics have eigenvalue
  $-\kappa(1 - F'(\pi^\star))$.  The fixed point is stable iff $F'(\pi^\star) < 1$,
  i.e., $G < 1$.  By the geometry of the S-curve, $\pi^\star_{\mathrm{low}}$
  and $\pi^\star_{\mathrm{high}}$ satisfy $F' < 1$ (stable), while
  $\pi^\star_{\mathrm{mid}}$ satisfies $F' > 1$ (unstable).
\end{proof}

\begin{remark}[Economic interpretation of the three equilibria]
  \label{rem:three_equil_econ}
  \leavevmode
  \begin{itemize}
    \item \textbf{Flywheel equilibrium} ($\pi^\star_{\mathrm{high}}$):
      High premium $\to$ aggressive issuance $\to$ rapid BTC-per-share
      accretion $\to$ justified high premium.  Self-reinforcing virtuous
      cycle.
    \item \textbf{Tipping point} ($\pi^\star_{\mathrm{mid}}$):
      Unstable separator.  Above it, the system is attracted to the flywheel;
      below it, to the low equilibrium.  Corresponds to the critical premium
      level at which the accumulation engine ``ignites.''
    \item \textbf{Low/spiral equilibrium} ($\pi^\star_{\mathrm{low}}$):
      Low premium $\to$ minimal issuance $\to$ negligible accretion $\to$
      justified low premium.  In the limit, if the premium turns negative,
      this becomes the death spiral (Section~\ref{sec:death_spiral}).
  \end{itemize}
\end{remark}

%-----------------------------------------------------------------------------
\subsubsection{Hysteresis}
\label{sec:hysteresis}
%-----------------------------------------------------------------------------

\begin{proposition}[Hysteresis in premium transitions]
  \label{prop:hysteresis}
  Consider a parameter path where $\Gamma$ increases slowly through the
  bifurcation point $\Gamma^*$ (e.g., due to rising $\alpha_N$ or falling $r$)
  and then decreases.
  %
  \begin{enumerate}
    \item As $\Gamma$ increases past $\Gamma^*$, the system (if starting at
      $\pi^\star_{\mathrm{low}}$) does \emph{not} jump to
      $\pi^\star_{\mathrm{high}}$; it remains at the low equilibrium until
      a sufficiently large shock pushes $\pi_t$ above the tipping point.
    \item If the system is at $\pi^\star_{\mathrm{high}}$ and $\Gamma$
      decreases below $\Gamma^*$, the flywheel equilibrium vanishes via
      saddle-node bifurcation, and the premium collapses discontinuously to
      the unique remaining (low) equilibrium.
    \item The path of equilibrium selection thus exhibits \emph{hysteresis}:
      the premium reached at a given $\Gamma$ depends on the direction of
      approach.
  \end{enumerate}
\end{proposition}

\begin{proof}
  Standard consequence of the saddle-node bifurcation structure.  When the
  stable branch (flywheel) and the unstable branch (tipping point) annihilate,
  no nearby fixed point exists, and the system transitions discontinuously
  to the only remaining attractor.
\end{proof}

%-----------------------------------------------------------------------------
\subsubsection{Death Spiral Mechanism}
\label{sec:death_spiral}
%-----------------------------------------------------------------------------

\begin{theorem}[Death spiral]
  \label{thm:death_spiral}
  Suppose the premium drops to $\pi_t < 0$ (discount to NAV).  Then:
  %
  \begin{enumerate}
    \item The accumulation option value becomes negligible:
      $V_{\mathrm{acc}} \approx 0$ since issuance at a discount is dilutive
      and the firm optimally refrains (the option is out of the money).

    \item The fair premium collapses to $\pi^\star \approx 0$
      (or negative if agency/overhead costs are included).

    \item If the firm has outstanding debt obligations requiring refinancing,
      it faces a dilemma: issue equity at a discount (dilutive to existing
      holders) or risk insolvency.

    \item The feedback loop reverses:
      \begin{equation}
        \label{eq:death_spiral_loop}
        \pi_t < 0 \;\Longrightarrow\;
        \alpha_N^{\mathrm{eff}} = 0 \;\Longrightarrow\;
        V_{\mathrm{acc}} \downarrow \;\Longrightarrow\;
        \pi^\star \downarrow \;\Longrightarrow\;
        \pi_t \downarrow\;\cdots
      \end{equation}

    \item A death spiral occurs when the premium discount exceeds the BTC
      growth expectation:
      \begin{equation}
        \label{eq:spiral_condition}
        e^{\pi_t} < \frac{D_t}{H_t S_t}
        \quad\Longleftrightarrow\quad
        E_t < D_t,
      \end{equation}
      meaning the market values equity below the debt burden.
  \end{enumerate}
\end{theorem}

\begin{proof}
  (1)--(2): When $\pi_t < 0$, the accretion rate
  $g(\pi, \Phi) = \alpha_N(e^\pi - 1)^+ \cdot L = 0$ since
  $(e^\pi - 1)^+ = 0$ for $\pi \le 0$.  Hence
  $V_{\mathrm{acc}} = 0$ and $\pi^\star = \log(1 + 0) = 0$.

  (3): With $\pi_t < 0$ and $\pi^\star \approx 0$, the restoring drift is
  $\kappa(0 - \pi_t) = -\kappa\pi_t > 0$, so the premium tends to recover.
  However, if the firm must issue equity to service debt, the dilutive
  issuance pushes $B_t$ down, reducing fundamental value and potentially
  deepening the discount.

  (4)--(5): The vicious cycle is triggered when forced dilutive issuance
  (to meet debt obligations) exceeds the natural mean-reversion, which
  occurs when $E_t < D_t$, i.e., the equity is insufficient to cover
  outstanding debt.
\end{proof}

%=============================================================================
\subsection{Mispricing Metric}
\label{sec:mispricing}
%=============================================================================

Given the endogenous fair premium $\pi^\star(\Phi_t)$, we can decompose the
observed premium into a ``justified'' and ``mispricing'' component.

\begin{definition}[Premium mispricing]
  \label{def:mispricing}
  The \emph{mispricing} (or \emph{premium gap}) is
  \begin{equation}
    \label{eq:delta_def}
    \Delta_t = \pi_t - \pi^\star(\Phi_t).
  \end{equation}
  Positive $\Delta_t$ indicates overvaluation relative to fundamentals;
  negative $\Delta_t$ indicates undervaluation.
\end{definition}

\begin{proposition}[Dynamics of the mispricing]
  \label{prop:delta_dynamics}
  Under the endogenous premium system~\eqref{eq:premium_endo} and assuming
  $\pi^\star(\Phi_t)$ varies slowly (Assumption~\ref{ass:steady_state}), the
  mispricing $\Delta_t$ satisfies
  \begin{equation}
    \label{eq:delta_ou}
    d\Delta_t \approx -\kappa_\Delta\,\Delta_t\,dt
      + \sigma_\pi\,dW_t^\pi
      - d\pi^\star(\Phi_t),
  \end{equation}
  where the \emph{mispricing mean-reversion speed} is
  \begin{equation}
    \label{eq:kappa_delta}
    \boxed{
      \kappa_\Delta = \kappa > \kappa_{\mathrm{eff}}.
    }
  \end{equation}
\end{proposition}

\begin{proof}
  Write $\pi_t = \pi^\star(\Phi_t) + \Delta_t$.  From~\eqref{eq:premium_endo}:
  \[
    d\pi_t = \kappa(\pi^\star - \pi_t)\,dt + \sigma_\pi\,dW_t^\pi
           = -\kappa\,\Delta_t\,dt + \sigma_\pi\,dW_t^\pi.
  \]
  Also, $d\pi_t = d\pi^\star + d\Delta_t$.  Subtracting:
  \[
    d\Delta_t = d\pi_t - d\pi^\star
              = -\kappa\,\Delta_t\,dt + \sigma_\pi\,dW_t^\pi - d\pi^\star.
  \]
  Under the slow-variation assumption, $d\pi^\star$ is small (driven by slow
  changes in $\Phi_t$), so the dominant dynamics of $\Delta_t$ are OU with
  speed $\kappa$.

  The comparison $\kappa_\Delta = \kappa > \kappa(1-G) = \kappa_{\mathrm{eff}}$
  follows since $G > 0$.
\end{proof}

\begin{remark}
  The mispricing reverts faster than the premium itself.  Intuitively,
  the premium has an ``extra'' source of persistence from the feedback loop
  (changes in $\pi$ shift $\pi^\star$ in the same direction, slowing
  reversion).  The mispricing, measured relative to the moving target
  $\pi^\star$, is free of this feedback and reverts at the ``bare'' speed
  $\kappa$.
\end{remark}

%-----------------------------------------------------------------------------
\subsubsection{Z-Score Bands}
\label{sec:zscore}
%-----------------------------------------------------------------------------

\begin{definition}[Mispricing Z-score]
  \label{def:zscore}
  The \emph{mispricing Z-score} is
  \begin{equation}
    \label{eq:zscore}
    Z_t^{\Delta} = \frac{\Delta_t}{\sigma_\Delta^{\mathrm{ss}}},
    \qquad
    \sigma_\Delta^{\mathrm{ss}}
    = \frac{\sigma_\pi}{\sqrt{2\kappa_\Delta}}
    = \frac{\sigma_\pi}{\sqrt{2\kappa}}.
  \end{equation}
\end{definition}

Under the OU approximation, $Z_t^\Delta$ is approximately standard normal in
the stationary distribution.  This yields the following classification:

\begin{center}
\begin{tabular}{lll}
  \toprule
  \textbf{Band} & \textbf{$|Z_t^\Delta|$} & \textbf{Interpretation} \\
  \midrule
  Fair value    & $< 1$   & Premium consistent with fundamentals \\
  Moderate      & $1$--$2$ & Mild over/undervaluation \\
  Extreme       & $> 2$   & Significant mispricing; mean-reversion expected \\
  \bottomrule
\end{tabular}
\end{center}

\begin{proposition}[Half-life of mispricing]
  \label{prop:halflife}
  The expected half-life of a mispricing deviation is
  \begin{equation}
    \label{eq:halflife}
    t_{1/2}^{\Delta} = \frac{\ln 2}{\kappa_\Delta} = \frac{\ln 2}{\kappa},
  \end{equation}
  which is shorter than the premium half-life
  $t_{1/2}^{\pi} = \ln 2 / \kappa_{\mathrm{eff}}$ by the factor $(1-G)$.
\end{proposition}

%=============================================================================
\subsection{Model Extensions and Discussion}
\label{sec:extensions}
%=============================================================================

%-----------------------------------------------------------------------------
\subsubsection{Market Impact and BTC Price Feedback}
\label{sec:market_impact}
%-----------------------------------------------------------------------------

Equation~\eqref{eq:btc_impact} includes a market-impact term
$\eta\,dH_t/H_{t^-}$ in the BTC price dynamics.  This introduces an
additional feedback channel:
\begin{equation}
  \label{eq:full_feedback}
  \pi_t \uparrow \;\to\;
  H_t \uparrow \;\to\;
  S_t \uparrow \text{ (via $\eta$)} \;\to\;
  \mathrm{NAV}_t \uparrow \;\to\;
  E_t \uparrow \;\to\;
  \text{issuance capacity} \uparrow \;\to\;
  H_t \uparrow\;\cdots
\end{equation}
The market-impact channel amplifies the feedback gain:

\begin{proposition}[Augmented feedback gain with market impact]
  \label{prop:augmented_gain}
  With market impact $\eta > 0$, the effective feedback gain becomes
  \begin{equation}
    \label{eq:augmented_gain}
    G_{\mathrm{aug}} = G + \eta\,\frac{\alpha_N L}{\kappa}\,\beta_{LS}(t),
  \end{equation}
  where $\beta_{LS}(t) = H_t S_t / (H_t S_t - D_t)$ is the structural
  leverage elasticity from the baseline model.  The second term captures the
  additional premium support from BTC price appreciation induced by MSTR
  purchases.
\end{proposition}

%-----------------------------------------------------------------------------
\subsubsection{Regime-Dependent Issuance}
\label{sec:regime_issuance}
%-----------------------------------------------------------------------------

The baseline model assumes a constant issuance intensity $\alpha_N$.  In
practice, the firm's issuance behavior is regime-dependent:
\begin{equation}
  \label{eq:regime_issuance}
  \alpha_N(\pi_t) =
  \begin{cases}
    \bar{\alpha}_N & \text{if } \pi_t > \pi_{\min} \text{ (premium regime)}, \\
    0              & \text{if } \pi_t \le \pi_{\min} \text{ (discount regime)},
  \end{cases}
\end{equation}
where $\pi_{\min} \ge 0$ is a threshold below which the firm suspends
issuance.  This can be smoothed via
$\alpha_N(\pi_t) = \bar{\alpha}_N \cdot \Phi((\pi_t - \pi_{\min})/\epsilon)$
for numerical stability.

%-----------------------------------------------------------------------------
\subsubsection{Connection to the Baseline Model}
\label{sec:connection}
%-----------------------------------------------------------------------------

The endogenous premium theory provides an economic foundation for the
parameters of the baseline OU model:
\begin{itemize}
  \item The long-run mean $\theta$ of the exogenous model is identified with
        the fair premium $\pi^\star(\Phi_t)$.  When state variables change
        slowly, $\theta \approx \pi^\star$ provides a good approximation.
  \item The effective mean-reversion speed is $\kappa_{\mathrm{eff}} = \kappa(1-G)$
        rather than the ``bare'' $\kappa$.  Empirical calibration of the OU
        model recovers $\kappa_{\mathrm{eff}}$, from which one can back out $G$
        if $\kappa$ is known.
  \item The IBGR metric from the baseline model is directly connected to
        $\alpha_N$: higher IBGR implies higher $\alpha_N$ and hence higher
        $\pi^\star$.
  \item The PMRI is reinterpreted: using the endogenous model,
        \begin{equation}
          \label{eq:pmri_endo}
          \mathrm{PMRI}_t^{\mathrm{endo}}
          = \frac{\pi_t - \pi^\star(\Phi_t)}{\sigma_\pi/\sqrt{2\kappa}}
          = Z_t^\Delta,
        \end{equation}
        which is the mispricing Z-score.  This is a strictly more informative
        metric than the exogenous PMRI, as it accounts for time-varying
        fundamentals.
\end{itemize}

%=============================================================================
\subsection{Summary of Main Results}
\label{sec:theory_summary}
%=============================================================================

\begin{enumerate}
  \item The NAV premium is the market value of a \emph{perpetual accumulation
    option}: the right to issue equity at a premium and buy BTC, increasing
    BTC-per-share (Proposition~\ref{prop:fair_premium}).

  \item The fair premium $\pi^\star$ satisfies a fixed-point equation with a
    closed-form steady-state approximation linking it to BTC volatility,
    issuance intensity, leverage, and discount rates
    (Proposition~\ref{prop:steady_state}).

  \item The reflexive feedback loop creates multiple equilibria (flywheel,
    tipping point, spiral) connected by saddle-node bifurcations with
    hysteresis (Theorem~\ref{thm:three_equilibria},
    Proposition~\ref{prop:hysteresis}).

  \item The mispricing $\Delta_t = \pi_t - \pi^\star$ reverts faster than the
    premium itself ($\kappa_\Delta = \kappa > \kappa_{\mathrm{eff}}$), enabling
    a more precise over/undervaluation metric
    (Proposition~\ref{prop:delta_dynamics}).

  \item The framework nests the baseline exogenous OU model as a special case
    and provides economic content to all its parameters.
\end{enumerate}
